% !Mode:: "TeX:UTF-8"

\section{项目简介}

随着嵌入式系统技术的快速发展和微控制器性能的不断提升，智能家电控制已成为现代生活的重要发展趋势。传统电风扇功能单一、操控不便，难以满足用户对舒适性和便捷性的更高要求。将嵌入式微控制器应用于家电控制领域，能够有效提升产品的智能化水平和用户体验。

STM32系列微控制器以其高性能、丰富的外设资源（如GPIO、定时器/PWM、ADC、通信接口等）、低功耗特性以及成熟的生态系统，成为嵌入式系统开发的理想选择。本项目依托嵌入式系统课程，旨在设计并实现一个基于STM32微控制器的遥控电风扇系统，将嵌入式软硬件设计技术应用于实际生活场景。

本设计的主要功能包括：
\begin{enumerate}
    \item \textbf{核心控制功能：} 利用STM32的定时器产生高精度PWM信号，实现对直流电机（风扇）的无级或多档位调速控制。
    \item \textbf{用户交互：} 实现红外遥控信号的接收与解码，使用户能够通过遥控器方便地开关风扇、调节风速档位。
    \item \textbf{增强功能：} 设计并实现定时关闭功能，满足用户预设使用时间的需求。
    \item \textbf{系统集成：} 完成从红外遥控接收、信号处理、核心控制逻辑（STM32软件）到电机驱动电路（硬件）的完整系统设计与集成。
    \item \textbf{实践验证：} 在真实的STM32开发板及外设搭建的硬件平台上进行系统测试与功能验证，确保设计的正确性和实用性。
\end{enumerate}

本项目不仅涵盖了嵌入式软件开发（C/C++）、硬件接口驱动（如红外接收、PWM输出）、控制算法（调速）等核心知识点，还涉及硬件电路设计（如电机驱动、电源管理）与系统调试等实践技能，是对嵌入式系统课程所学知识的综合应用与实践。

% 本报告后续章节将详细阐述系统设计的具体方案。第\ref{sec:hardware}章介绍系统整体架构与各模块的硬件设计；第\ref{sec:software}章详细描述基于STM32的软件设计与关键代码实现；第\ref{sec:testing}章展示系统测试方法与最终实现效果；最后，第\ref{sec:conclusion}章对项目进行总结并探讨可能的改进方向。